\exercisesection{2}

\begin{enumerate}
\def\labelenumi{\arabic{enumi}.}
\tightlist
\item
  \begin{enumerate}
  \def\labelenumii{(\alph{enumii})}
  \tightlist
  \item
    Show that in the Noetherian polynomial ring \(\mathbf{A} = \mathbf{Z}[\mathbf{X}]\) in one indeterminate over \(\mathbf{Z}\) , the ideal \(m = 2A + AX\) is maximal and that the ideal \(q = 4A + AX\) is m-primary, but is not equal to a power of \(m\) .
  \end{enumerate}
\end{enumerate}

\begin{enumerate}
\def\labelenumi{(\alph{enumi})}
\setcounter{enumi}{1}
\item
  In the ring \(\mathbf{B} = \mathbf{Z}[2\mathbf{X},\mathbf{X}^2,\mathbf{X}^3] \subset \mathbf{A}\) which is Noetherian, show that the ideal \(\mathfrak{p} = 2\mathrm{BX} + \mathrm{BX}^2\) is prime, but that \(\mathfrak{p}^2\) is not \(\mathfrak{p}\) -primary, although its radical is equal to \(\mathfrak{p}\) .
\item
  Let \(\mathbf{K}\) be a commutative field and \(\mathbf{A}\) the quotient ring of \(\mathbf{K}[\mathbf{X},\mathbf{Y},\mathbf{Z}]\) by the ideal generated by \(Z^2 -\mathbf{XY}\) ; let \(x,y,z\) be the canonical images of \(\mathbf{X},\mathbf{Y},\mathbf{Z}\) in \(\mathbf{A}\) . Show that \(\mathfrak{p} = \mathbf{A}x + \mathbf{A}z\) is prime, that \(\mathfrak{p}^2\) is not primary and that \(\mathfrak{p}^2 = \mathfrak{a}\cap \mathfrak{b}^2\) is a primary decomposition of \(\mathfrak{p}^2\) , where \(\mathfrak{a} = \mathbf{A}x\) and \(\mathfrak{b} = \mathbf{A}x + \mathbf{A}y + \mathbf{A}z\) .
\end{enumerate}

\begin{enumerate}
\def\labelenumi{\arabic{enumi}.}
\setcounter{enumi}{1}
\item
  In the example of Exercise 11 (b) of §1, show that \(\mathfrak{a} = \mathfrak{m}^2 \cap \mathfrak{p}_1 \cap \mathfrak{p}_2\) is a reduced primary decomposition of \(\mathfrak{a}\) in \(\mathbf{A}\) and that the saturation of \(\mathfrak{a}\) with respect to \(\mathfrak{m}\) is equal to \(\mathfrak{a}\) (and is therefore not primary).
\item
  Let A be a Noetherian ring and M a finitely generated A-module. Let Q be a p-primary submodule in M. The greatest lower bound of the integers \(n \geqslant 1\) such that \(p^{n}M \subset Q\) is called the exponent of Q in M and denoted by \(e(M/Q)\) .
\end{enumerate}

Let \((Q_{\lambda})_{\lambda \in L}\) be a family of \(p\) -primary submodules in M. For \(Q = \bigcap_{\lambda \in L} Q_{\lambda}\) to be p-primary in M, it is necessary and sufficient that the family \((e(\mathbf{M}/\mathbf{Q}_{\lambda}))_{\lambda\in\mathbf{L}}\) be bounded above; then \(e(\mathbf{M}/\mathbf{Q})\geqslant e(\mathbf{M}/\mathbf{Q}_{\lambda})\) for all \(\lambda\in L\) .

¶ 4. Let A be a Noetherian ring and M a finitely generated A-module. Let p be an element of Ass(M) and let \(\mathfrak{G}_{\mathfrak{p}}\) be the set of submodules Q of M such that \(\operatorname{Ass}(M/Q) = \{\mathfrak{p}\}\) and \(\operatorname{Ass}(Q) = \operatorname{Ass}(M) - \{\mathfrak{p}\}\) (a set which is not empty by § 1, no. 1, Proposition 4). We write \(e_{\mathfrak{p}}(M) = \inf_{\mathbf{Q} \in \mathfrak{G}_{\mathfrak{p}}} e(M/Q)\) (Exercise 3).

\begin{enumerate}
\def\labelenumi{(\alph{enumi})}
\item
  Let \(n_{\mathfrak{p}}\) be an integer \(\geqslant e_{\mathfrak{p}}(\mathbf{M})\) and let \(\mathfrak{F}(n_{\mathfrak{p}})\) be the subset of \(\mathfrak{G}_{\mathfrak{p}}\) consisting of the submodules \(\mathbf{Q}\) such that \(e(\mathbf{M} / \mathbf{Q}) \leqslant n_{\mathfrak{p}}\) . Show that \(\mathfrak{F}(n_{\mathfrak{p}})\) has a least element \(\mathbf{Q}(\mathfrak{p}, n_{\mathfrak{p}})\) .
\item
  Let \((n_{\mathfrak{p}})_{\mathfrak{p} \in \operatorname{Ass}(\mathbf{M})}\) be a family of integers such that \(n_{\mathfrak{p}} \geqslant e_{\mathfrak{p}}(\mathbf{M})\) for all \(\mathfrak{p} \in \operatorname{Ass}(\mathbf{M})\) . Show that the submodules \(\mathbf{Q}(\mathfrak{p}, n_{\mathfrak{p}})\) corresponding to this family form a reduced primary decomposition of \(\{0\}\) in \(\mathbf{M}\) said to be canonically determined by the family \((n_{\mathfrak{p}})\) . If we take \(n_{\mathfrak{p}} = e_{\mathfrak{p}}(\mathbf{M})\) for all \(\mathfrak{p} \in \operatorname{Ass}(\mathbf{M})\) , the primary decomposition consisting of the \(\mathbf{Q}(\mathfrak{p}) = \mathbf{Q}(\mathfrak{p}, e_{\mathfrak{p}}(\mathbf{M}))\) is called the canonical primary decomposition of \(\{0\}\) in \(\mathbf{M}\) . Let \(\{0\} = \bigcap_{\mathfrak{p} \in \operatorname{Ass}(\mathbf{M})} \mathbf{Q}'(\mathfrak{p})\) be any reduced primary decomposition of \(\{0\}\) in \(\mathbf{M}\) . Show that \(e(\mathbf{M}/\mathbf{Q}(\mathfrak{p})) \leqslant e(\mathbf{M}/\mathbf{Q}'(\mathfrak{p}))\) for all \(\mathfrak{p} \in \operatorname{Ass}(\mathbf{M})\) and that, if \(e(\mathbf{M}/\mathbf{Q}(\mathfrak{p})) = e(\mathbf{M}/\mathbf{Q}'(\mathfrak{p}))\) for some \(\mathfrak{p}\) , then \(\mathbf{Q}(\mathfrak{p}) \subset \mathbf{Q}'(\mathfrak{p})\) (``Theorem of Ortiz'').
\item
  Show that \(\mathbf{Q}(\mathfrak{p}, n_{\mathfrak{p}})\) is the saturation of \(\mathfrak{p}^{n_{\mathfrak{p}}} \mathbf{M}\) with respect to \(\mathfrak{p}\) (use § 1, no. 2, Proposition 6); \(\mathfrak{p}\) is the least element of \(\operatorname{Supp}(\mathbf{M} / \mathfrak{p}^n\mathbf{M})\) .
\item
  Let S be a multiplicative subset of A not meeting a prime ideal \(\mathfrak{p} \in \operatorname{Ass}(M)\) and let Q be a p-primary submodule of M. Show that \(e(\mathrm{M}/\mathrm{Q}) = e(\mathrm{S}^{-1}\mathrm{M}/\mathrm{S}^{-1}\mathrm{Q})\) . Show that, if S is a multiplicative subset of A, \(\Phi\) the set of \(\mathfrak{p} \in \operatorname{Ass}(M)\) such that \(S \cap \mathfrak{p} = \varnothing\) and N the saturation of \(\{0\}\) in M with respect to S, the \(\mathrm{Q}(\mathfrak{p}, n_{\mathfrak{p}})\) where \(\mathfrak{p} \in \Phi\) form the reduced primary decomposition of N in M canonically determined by the family \((n_{\mathfrak{p}})_{\mathfrak{p} \in \Phi}\) and the \(S^{-1}\mathrm{Q}(\mathfrak{p}, n_{\mathfrak{p}})\) where \(\mathfrak{p} \in \Phi\) the reduced primary decomposition of \(\{0\}\) in \(S^{-1}M\) canonically determined by the family \((n_{\mathfrak{p}})_{\mathfrak{p} \in \Phi}\) . Consider the particular case of canonical primary decompositions.
\end{enumerate}

\begin{enumerate}
\def\labelenumi{\arabic{enumi}.}
\setcounter{enumi}{4}
\item
  Let L be a finitely generated free Z-module and T a finite commutative group whose order is a power of a prime number p.~Show that the canonical primary decomposition (Exercise 4) of \(\{0\}\) in \(M = L \oplus T\) is \(\{0\} = p^{n}L \cap T\) , where n is the least integer \(\geqslant 0\) such that \(p^{n}T = 0\) (note that \(p^{n}M = p^{n}L\) ).
\item
  In Exercise 5 of § 1, \{0\} is primary in the A-module a relative to the prime ideal \{0\} of A. Show that the submodule AX of a is primary with respect to a relative to a prime ideal \(\neq\{0\}\) and that the submodule AXY of a is not primary with respect to a.
\item
  Let \(\mathbf{A}\) be a Noetherian ring, \(\mathbf{M}\) a finitely generated A-module and \((\mathfrak{p}_i)_{1\leqslant i\leqslant n}\) a sequence obtained by arranging the elements of \(\operatorname {Ass}(\mathbf{M})\) in any order. Show that there exists a composition series \((\mathbf{M}_i)_{0\leqslant i\leqslant n}\) of \(\mathbf{M}\) such that, for \(0\leqslant i\leqslant n - 1\) , \(\mathbf{M}_{i + 1}\) is \(p_i\) primary in \(\mathbf{M}_i\) (use Proposition 4 of no. 3).
\end{enumerate}

¶ 8. Let A be a Noetherian ring, E a finitely generated A-module, F a submodule of E and m an ideal of A. Let \(\mathbf{F} = \bigcap_{\mathfrak{p} \in \operatorname{Ass}(\mathbf{E}/\mathbf{F})} \mathbf{Q}(\mathfrak{p})\) be a reduced primary decomposition of F in E. Show that the closure of F in E under the m-adic topology on E is equal to \(\bigcap_{\mathfrak{p} \in \Phi} \mathbf{Q}(\mathfrak{p})\) , where \(\Phi\) is the set of \(\mathfrak{p} \in \operatorname{Ass}(\mathbf{E}/\mathbf{F})\) such that \(\mathfrak{p} + \mathfrak{m} \neq \mathbf{A}\) . (Consider first the case where F is p-primary in E and show that, if \(\mathfrak{p} + \mathfrak{m} = \mathbf{A}\) , F is dense in E; proceed to the general case using Exercise 15 of § 1 and Chapter III, § 3, no. 4, Corollary 1 to Theorem 3).

\begin{enumerate}
\def\labelenumi{\arabic{enumi}.}
\setcounter{enumi}{8}
\item
  Let A be a Noetherian ring, m a maximal ideal of A and M an A-module such that \(\{0\}\) is m-primary in M. If \(\hat{A}\) is the Hausdorff completion of A with respect to the m-adic topology, show that the canonical mapping \(M \to M \otimes_{A} \hat{A}\) is bijective (consider first the case where M is finitely generated, using Proposition 7 of no. 5; in the general case, consider M as the direct limit of its finitely generated submodules).
\item
  Let A be a Noetherian ring and M an A-module. Show that the following properties are equivalent:
\end{enumerate}

\((\alpha)\) Every finitely generated submodule of \(\mathbf{M}\) is of finite length.

(β) M is a direct limit of A-modules of finite length.

(γ) Every element \(\mathfrak{p} \in \operatorname{Ass}(\mathbf{M})\) is a maximal ideal of A.

(δ) Every element \(\mathfrak{p} \in \operatorname{Supp}(\mathbf{M})\) is a maximal ideal of A.

\begin{enumerate}
\def\labelenumi{\arabic{enumi}.}
\setcounter{enumi}{10}
\tightlist
\item
  Let A be a ring, M an A-module and N a submodule of M. The set of \(a \in A\) such that the homothety of ratio a on M/N is almost nilpotent is called the radical of N in M and denoted by \(\mathfrak{r}_{\mathrm{M}}(\mathrm{N})\) . Show that \(\mathfrak{r}_{\mathrm{M}}(\mathrm{N})\) is an ideal of A and that, for all \(p \in \operatorname{Ass}_{f}(\mathrm{M}/\mathrm{N})\) , \(\mathfrak{r}_{\mathrm{M}}(\mathrm{N}) \subset p\) . If \(N_{1}, N_{2}\) are two submodules of M, then \(\mathfrak{r}_{\mathrm{M}}(\mathrm{N}_{1} \cap \mathrm{N}_{2}) = \mathfrak{r}_{\mathrm{M}}(\mathrm{N}_{1}) \cap \mathfrak{r}_{\mathrm{M}}(\mathrm{N}_{2})\) ; if a is an ideal of A, \(\mathfrak{r}_{\mathrm{M}}(\mathfrak{a}\mathrm{N}) \supset \mathfrak{r}(\mathfrak{a}) \cap \mathfrak{r}_{\mathrm{M}}(\mathrm{N})\) and \(\mathfrak{r}_{\mathrm{A}}(\mathfrak{a}) = \mathfrak{r}(\mathfrak{a})\) .
\end{enumerate}

¶ 12. Let A be a ring, M an A-module and N a submodule of M.

\begin{enumerate}
\def\labelenumi{(\alph{enumi})}
\item
  For \(\operatorname{Ass}_f(\mathbf{M} / \mathbf{N})\) (§ 1, Exercise 17) to be reduced to a single element \(\mathfrak{p}\) , it is necessary and sufficient that \(\mathbf{N} \neq \mathbf{M}\) and that, for all \(a \in \mathbf{A}\) , the homothety with ratio \(a\) on \(\mathbf{M} / \mathbf{N}\) be injective or almost nilpotent; then \(\mathfrak{p} = \mathfrak{r}_{\mathbf{M}}(\mathbf{N})\) (Exercise 11) and we also say that \(\mathbf{N}\) is primary (or \(\mathfrak{p}\) -primary) in \(\mathbf{M}\) . For every submodule \(\mathbf{M}'\) of \(\mathbf{M}\) such that \(\mathbf{N} \subset \mathbf{M}'\) and \(\mathbf{N} \neq \mathbf{M}'\) , \(\mathbf{N}\) is then \(\mathfrak{p}\) -primary in \(\mathbf{M}'\) .
\item
  If \(\mathfrak{r}_{\mathbb{M}}(\mathbf{N})\) is a maximal ideal of \(\mathbf{A}\) , show that \(\mathbf{N}\) is primary in \(\mathbf{M}\) . In particular, every ideal \(q\) of \(\mathbf{A}\) which is only contained in a single prime ideal \(m\) (necessarily maximal) is \(m\) -primary.
\item
  Let \(\mathbf{A}\) be an integral domain and \(x\) an element of \(\mathbf{A}\) such that \(\mathfrak{p} = \mathbf{A}x\) is prime; then, for every integer \(n \geqslant 1\) , \(\mathrm{Ax}^n\) is \(\mathfrak{p}\) -primary. Show that this result does not extend to the case where A is not an integral domain (take A = B/b, where B = K{[}X, Y{]} (K being a field), b = BX² + BXY).
\item
  If A is an absolutely flat ring (Chapter I, § 2, Exercise 17), the primary ideals of A are identical with the prime ideals of A.
\item
  Give an example of a Z-module M such that \(\{0\}\) is p-primary (for a prime number p) but that for no \(a \neq 0\) in Z is the homothety \(a_{M}\) nilpotent (cf.~Chapter II, §2, Exercise 3).
\end{enumerate}

\begin{enumerate}
\def\labelenumi{\arabic{enumi}.}
\setcounter{enumi}{12}
\tightlist
\item
  \begin{enumerate}
  \def\labelenumii{(\alph{enumii})}
  \tightlist
  \item
    Let \(u: \mathbf{M} \to \mathbf{M}'\) be a surjective A-module homomorphism. Show that, if \(\mathbf{N}'\) is a primary submodule in \(\mathbf{M}'\) , \(\mathbf{N} = \bar{u}^{-1}(\mathbf{N}')\) is primary in \(\mathbf{M}\) and \(\mathfrak{r}_{\mathbf{M}}(\mathbf{N}) = \mathfrak{r}_{\mathbf{M}'}(\mathbf{N}')\) .
  \end{enumerate}
\end{enumerate}

\begin{enumerate}
\def\labelenumi{(\alph{enumi})}
\setcounter{enumi}{1}
\item
  Every intersection of a non-empty finite family of p-primary submodules in an A-module M is p-primary in M. Does the proposition extend to arbitrary intersections?
\item
  For an A-module M to be such that \(\{0\}\) is p-primary in M, it is necessary and sufficient that, for every submodule \(N \neq 0\) of M, \(\operatorname{Supp}(N) = V(p)\) . (Note that, for all \(x \in M\) , M contains a submodule isomorphic to \(A/\operatorname{Ann}(x)\) .)
\end{enumerate}

\begin{enumerate}
\def\labelenumi{\arabic{enumi}.}
\setcounter{enumi}{13}
\tightlist
\item
  \begin{enumerate}
  \def\labelenumii{(\alph{enumii})}
  \tightlist
  \item
    Generalize Proposition 3 of no. 1 to arbitrary rings.
  \end{enumerate}
\end{enumerate}

\begin{enumerate}
\def\labelenumi{(\alph{enumi})}
\setcounter{enumi}{1}
\tightlist
\item
  Let M be an A-module such that, for every multiplicative subset S of A, the kernel of the canonical mapping \(M \rightarrow S^{-1}M\) is \(\{0\}\) or M. Show that \(\{0\}\) is primary in M (consider for all \(a \neq 0\) in A the multiplicative subset of the \(a^{n}\) ( \(n \geqslant 0\) ).
\end{enumerate}

¶ 15. Let q be a primary ideal in a ring A and p its radical. Show that in the polynomial ring B = A{[}X{]} the ideal Bq is primary and has radical the ideal Bp. (Let \(f(X) = \sum_{j} a_j X^j\) and \(g(X) = \sum_{j} b_j X^j\) be two non-constant polynomials such that \(fg \in Bq\) and \(f \notin Bp\) . Let \(a_m\) be the coefficient of smallest index not belonging to p; let a be the ideal of A generated by \(a_0, a_1, \ldots, a_{m-1}\) ; there exists an integer k such that \(a^k \subset q\) . Let \(q_i\) be the transporter of \(a^{k-i}\) in q for \(i \leqslant k\) . Show by induction on i that \(g \in Bq_i\) , arguing by reductio ad absurdum.

\begin{enumerate}
\def\labelenumi{\arabic{enumi}.}
\setcounter{enumi}{15}
\item
  In a ring A let p be a non-maximal prime ideal and q a p-primary ideal distinct from p.~Let x be an element of p - q and y an element of A - p such that \(p + Ay \neq A\) . Show that the ideal \(a = q + Axy\) , such that \(p \supset a \supset q\) , is not primary (note that \(x \in a\) is impossible).
\item
  Let M be an A-module and N a p-primary submodule in M.
\end{enumerate}

\begin{enumerate}
\def\labelenumi{(\alph{enumi})}
\item
  Let F be a submodule of M not contained in N. Show that the transporter N:F is an ideal of A contained in p.
\item
  Let \(\mathfrak{a}\) be an ideal of \(\mathbf{A}\) . Show that if \(\mathfrak{a} \subset \mathbb{N}: \mathbb{M}\) , then \(\mathbb{N}: a = \mathbb{M}\) . If \(\mathfrak{a} \notin \mathbb{N}: \mathbb{M}\) , \(\mathbb{N}: \mathfrak{a}\) is a \(\mathfrak{p}\) -primary submodule in \(\mathbb{M}\) . If \(\mathfrak{a} \notin \mathfrak{p}\) , then \(\mathbb{N}: \mathfrak{a} = \mathbb{N}\) .
\end{enumerate}

\begin{enumerate}
\def\labelenumi{\arabic{enumi}.}
\setcounter{enumi}{17}
\tightlist
\item
  \begin{enumerate}
  \def\labelenumii{(\alph{enumii})}
  \tightlist
  \item
    Let \(\mathbf{A}\) be a ring and \(\mathbf{M}\) an \(\mathbf{A}\) -module. Show that, if \(\mathfrak{p}\) is a minimal element of \(\operatorname{Ass}_f(\mathbf{M})\) , the saturation of \(\{0\}\) with respect to \(\mathfrak{p}\) in \(\mathbf{M}\) is \(\mathfrak{p}\) -primary in \(\mathbf{M}\) (cf.~§ 1, Exercise 17 (b)).
  \end{enumerate}
\end{enumerate}

\begin{enumerate}
\def\labelenumi{(\alph{enumi})}
\setcounter{enumi}{1}
\item
  Let p be a prime ideal of A. Show that for every integer n \textgreater{} 0 the saturation in A of \(p^{n}\) with respect to p is p-primary in A (use Exercise 14 (a)); this saturation is denoted by \(p^{(n)}\) and is called the n-th symbolic power of p.
\item
  Show that, if every prime ideal of A is maximal, then, for every A-module M and every submodule N of M, N is the intersection of a family (finite or otherwise) of primary submodules in M (use (a) and Exercise 17 (i) of § 1).
\end{enumerate}

* 19. Determine the primary ideals in the ring of a valuation of height 2 (cf.~Chapter VI); deduce an example of a ring where there are ideals which are not intersections of a family (finite or otherwise) of primary ideals, although there exist in the ring only two distinct prime ideals. *

¶ 20. The notion of primary decomposition and that of reduced primary decomposition may be defined for any ring A as in nos. 2 and 3 (using the notion of primary submodule defined in Exercise 12).

\begin{enumerate}
\def\labelenumi{(\alph{enumi})}
\item
  Generalize Propositions 4 and 6 by replacing Ass by \(\mathbf{Ass}_f\) throughout.
\item
  Let E be a finitely generated A-module and F a submodule of E. For every multiplicative subset S of A let \(\operatorname{sat}_{\mathbf{S}}(\mathbf{F})\) denote the saturation of F in E with respect to S. Show that, if F admits a primary decomposition in E, every saturation of F in E is of the form F: (a) for some \(a \in A\) . (Let \(p_{i}\) ( \(1 \leqslant i \leqslant r\) ) be the elements of \(\operatorname{Ass}_{f}(\operatorname{E}/\operatorname{F})\) . Show first that, using (a), that for every multiplicative subset S of A there exists \(b \in A\) such that, if T is the multiplicative subset of the \(b^{n}\) ( \(n \geqslant 0\) ), then \(\operatorname{sat}_{\mathbf{T}}(\mathbf{F}) = \operatorname{sat}_{\mathbf{S}}(\mathbf{F})\) , taking b such that \(b \in p_{i}\) if \(p_{i} \cap S \neq 0\) and \(b \notin p_{i}\) otherwise. Then prove that we may take \(a = b^{m}\) for m sufficiently large.)
\item
  Give an example of a \(\mathbf{Z}\) -module E such that \(\{0\}\) is primary with respect to E but there exist saturations \(\mathrm{sat}_{\mathbb{S}}(0)\) which are not of the form 0: (a) (cf.~Algebra, Chapter VII, § 2, Exercise 3).
\end{enumerate}

\begin{enumerate}
\def\labelenumi{\arabic{enumi}.}
\setcounter{enumi}{20}
\tightlist
\item
  Let A be a ring, F a monic polynomial of A{[}X{]}, B the ring A{[}X{]}/F.A{[}X{]}, m a maximal ideal of A, k = A/m the quotient field and \(\bar{F}\) the canonical image of F in k{[}X{]}.
\end{enumerate}

\begin{enumerate}
\def\labelenumi{(\alph{enumi})}
\item
  Let \(\overline{\mathbf{F}} = \prod_{i=1}^{n} f_i^{e_i}\) , where the \(f_i\) are distinct irreducible polynomials of \(k[\mathbf{X}]\) . Let \(\mathbf{F}_i \in \mathbf{A}[\mathbf{X}]\) be a monic polynomial such that its canonical image in \(k[\mathbf{X}]\) is \(f_i\) . Let \(\mathfrak{M}_i\) denote the ideal \(\mathsf{mB} + \mathbf{F}_i\) . B of B; show that the \(\mathfrak{M}_i\) are distinct and are the only maximal ideals of B containing \(\mathsf{mB}\) (note that \(\mathbf{B}/\mathsf{mB}\) is an algebra of finite rank over \(\mathbf{A}/\mathsf{m} = k\) generated by the roots of \(\overline{\mathbf{F}}\) ).
\item
  For all i let us write \(Q_{i} = mB + F_{i}^{e_{i}}\) . B. Show that the \(Q_{i}\) are \(M_{i}\) -primary in B and that \(mB = Q_{1} \cap Q_{2} \cap \cdots \cap Q_{n} = Q_{1}Q_{2} \cdots Q_{n}\) .
\item
  For B to be a local ring, it is necessary and sufficient that A be a local ring and that \(\overline{F}\) be a power of an irreducible polynomial of k{[}X{]}.
\end{enumerate}

¶ 22. Let E be an A-module and F a submodule of E.

\begin{enumerate}
\def\labelenumi{(\alph{enumi})}
\item
  Let \(\mathfrak{p}\) be a prime ideal of \(\mathbf{A}\) and \(a\) an element of \(\mathbf{A}\) such that \(a \notin \mathfrak{p}\) and \(\mathbf{Q} = \mathbf{F}\) : (a) is \(\mathfrak{p}\) -primary in \(\mathbf{E}\) . Show that \(\mathbf{F} = \mathbf{Q} \cap (\mathbf{F} + a\mathbf{E})\) .
\item
  Suppose that there exist \(b \in \mathbf{A}\) and \(x \in \mathbf{E}\) such that \(b^n x \notin \mathbf{F}\) for all \(n > 0\) . Show that, if there exists an integer \(n > 0\) such that \(\mathbf{F}: b^{n+1} = \mathbf{F}: b^n\) , then \(\mathbf{F} = (\mathbf{F} + \mathbf{Ab}^n x) \cap (\mathbf{F}: (b^n))\) .
\item
  Suppose that \(\mathbf{F}\) is irreducible with respect to \(\mathbf{E}\) (Algebra, Chapter II, § 2, Exercise 16) and consider the three following properties:
\end{enumerate}

(α) F is primary in E.

(β) For every prime ideal p of A, the saturation of F with respect to p in E is of the form F: (a) for some \(a \notin p\) .

\((\gamma)\) For all \(b \in \mathbf{A}\) , the sequence \((\mathbf{F}:(b^n))_{n \geqslant 1}\) is stationary.

Show that each of conditions \((\beta)\) , \((\gamma)\) implies \((\alpha)\) (use (a) and (b) and Exercise 18 (a)). If E is finitely generated, the three conditions \((\alpha)\) , \((\beta)\) and \((\gamma)\) are equivalent (cf.~Exercise 20 (b) and 20 (c)).

\begin{enumerate}
\def\labelenumi{(\alph{enumi})}
\setcounter{enumi}{3}
\item
  Let \(\mathbf{K}\) be a field, \(\mathbf{A}\) the ring \(\mathbf{K}[\mathbf{X},\mathbf{Y}]\) and \(\mathfrak{m}\) the maximal ideal \(\mathbf{AX} + \mathbf{AY}\) of \(\mathbf{A}\) ; show that the ideal \(\mathfrak{m}^2\) is primary but not irreducible.
\item
  Let A be a not necessarily commutative ring and E a left Noetherian A-module. Show that every submodule F of E is the intersection of a finite family of irreducible submodules in E (consider a submodule of E which is maximal among those which are not finite intersections of irreducible submodules).
\item
  Deduce from (e) and (c) a new proof of Theorem 1 of no. 2.
\end{enumerate}

¶ 23. Let A be a ring. An A-module E is called Laskerian if it finitely generated and if every submodule of E has a primary decomposition in E.

\begin{enumerate}
\def\labelenumi{(\alph{enumi})}
\tightlist
\item
  For a finitely generated A-module E to be Laskerian, it is necessary and sufficient that it satisfy the two following axioms:
\end{enumerate}

\((\mathbf{LA}_{\mathbf{I}})\) For every submodule \(\mathbf{F}\) of \(\mathbf{E}\) and every prime ideal \(\mathfrak{p}\) of \(\mathbf{A}\) , the saturation of \(\mathbf{F}\) with respect to \(\mathfrak{p}\) in \(\mathbf{E}\) is of the form \(\mathbf{F}:(a)\) for some \(a\notin\mathfrak{p}\) .

\((\mathbf{L}\mathbf{A}_{\mathrm{II}})\) For every submodule \(\mathbf{F}\) of \(\mathbf{E}\) , every decreasing sequence \((\mathrm{sat}_{\mathbf{S}_{n}}(\mathbf{F}))\) (where \((\mathbf{S}_{n})\) is any decreasing sequence of multiplicative subsets of \(\mathbf{A}\) ) is stationary.

(To show that the conditions are sufficient, prove first, using \((\mathrm{LA}_{\mathrm{I}})\) and Exercises 18 (a) and 22 (a), that, for every submodule F of E, there exists a submodule Q of E which is primary for some ideal \(p \supset F: E\) and a submodule \(G = F + aE\) , where \(a \notin p\) , such that \(F = Q \cap G\) . Then argue by reductio ad absurdum: show that there would exist an infinite sequence \((\mathbf{Q}_{n})\) of \(p_{n}\) -primary submodules of E and a strictly increasing sequence \((\mathbf{G}_{n})\) of submodules of E such that, if we write \(H_{n} = Q_{1} \cap Q_{2} \cap \cdots \cap Q_{n}\) : (1) \(F = H_{n} \cap G_{n}\) ; (2) \(G_{n}\) is maximal among the submodules G containing \(G_{n-1}\) and such that \(\mathbf{F} = \mathbf{H}_n \cap \mathbf{G}\) ; (3) there exists \(a_n \notin p_n\) such that \(a_n \mathrm{E} \subset \mathrm{G}_n\) . Show then that, if \(\mathrm{S}_n = \bigcap_{1 \leqslant i \leqslant n} (\mathrm{A} - p_i)\) , \(\mathrm{S}_n\) meets \(\mathrm{G}_n: \mathrm{E}\) and deduce that \(\operatorname{sat}_{\mathrm{S}_n}(\mathrm{F}) = \mathrm{H}_n\) . Conclude with the aid of \((\mathrm{LA}_{\mathrm{II}})\) .

* (b) Give an example of a ring A such that the A-module A does not satisfy \((\mathrm{LA}_{\mathrm{I}})\) but where every ideal \(a \subset A\) is irreducible and the set of saturations \(\operatorname{sat}_{\mathrm{S}}(a)\) for all the multiplicative subsets S of A is finite (cf.~Exercise 19). *

\begin{enumerate}
\def\labelenumi{(\alph{enumi})}
\setcounter{enumi}{2}
\item
  Let K be a field, \(B = K[[X]]\) the algebra of formal power series over K and m the maximal ideal of B consisting of the formal power series without constant term. In the product algebra \(B^{N}\) , consider the subalgebra A generated by 1 and the ideal \(n = m^{(N)}\) . Show that in A the only prime ideals distinct from n are the ideals the direct sums in n of all the component ideals except one; every strictly increasing sequence of prime ideals of A then has at most two elements. Let \(a \neq A\) be an ideal of A and let S be a multiplicative subset of A not meeting a; if \(a_{n}\) (resp. \(S_{n}\) ) is the projection of a (resp. S) onto the n-th factor \(B_{n}\) of \(B^{N}\) , show that \(\text{sat}_{S}(a)\) is the direct sum of the ideals \(\text{sat}_{S_{n}}(a_{n})\) (observe that, if \(s \in S\) , \(s_{n} \in S_{n}\) is the n-th projection of s and \(x_{n} \in \text{sat}_{S_{n}}(a_{n})\) is such that \(s_{n}x_{n} \in a_{n}\) , then \(s^{2}x_{n} \in a\) ). Deduce that A satisfies axiom \((\mathbf{LA}_{\mathbf{I}})\) but not \((\mathbf{LA}_{\mathbf{II}})\) (use Exercise 20 (b)).
\item
  Show that, if an A-module E satisfies axiom \((\mathrm{LA}_{\mathrm{I}})\) , every submodule of E is the intersection of a family (finite or otherwise) of primary submodules in E (argue as in (a), by transfinite induction).
\item
  Let E be a finitely generated A-module and \(E' \subset E\) a finitely generated submodule. For E to be Laskerian, it is necessary and sufficient that \(E'\) and \(E/E'\) be so. In particular, if A is a Laskerian ring, every finitely generated A-module is Laskerian. If E is a Laskerian A-module, \(S^{-1}E\) is a Laskerian \(S^{-1}A\) -module for every multiplicative subset S of A.
\end{enumerate}

¶ 24. In a Laskerian ring A (Exercise 23), let \(p_1, p_2, p_3\) be three distinct prime ideals such that \(p_1 \subset p_2 \subset p_3\) ; let \(x_2 \in p_2 - p_1, x_3 \in p_3 - p_2\) . Replacing if need be \(p_2\) and \(p_3\) by prime ideals contained in \(p_2, p_3\) respectively and containing respectively \(x_2\) and \(x_3\) , we may suppose that \(p_2\) is minimal in \(\operatorname{Ass}_f(A / (p_1 + Ax_2))\) and \(p_3\) is minimal in \(\operatorname{Ass}_f(A / p_2 + Ax_3)\) . Consider the ideal \(a = p_1 + x_2p_3\) .

\begin{enumerate}
\def\labelenumi{(\alph{enumi})}
\item
  Show that \(x_{2} \notin \mathfrak{a}\) and \(x_{3}^{k} \notin \mathfrak{a}\) for every integer \(k > 0\) ; deduce that \(\mathfrak{a}\) is not primary in A. Show that \(\mathfrak{p}_{2}\) is a minimal element of \(\operatorname{Ass}_f(\mathbf{A} / \mathfrak{a})\) .
\item
  Let \(\mathfrak{a} = \bigcap_{i=1}^{n} \mathfrak{q}_i\) be a reduced primary decomposition of \(\mathfrak{a}\) in \(\mathbf{A}\) and let \(\mathfrak{p}_i'\) be the radical of \(\mathfrak{q}_i\) ; suppose that \(\mathfrak{p}_3 \notin \mathfrak{p}_i'\) for \(1 \leqslant i \leqslant s\) , \(\mathfrak{p}_3 \subset \mathfrak{p}_i'\) for \(s + 1 \leqslant i \leqslant n\) ; show that of necessity \(s < n\) and \(x_2 \in \mathfrak{q}_i\) for \(1 \leqslant i \leqslant s\) . Show that there exists an index \(i \geqslant s + 1\) such that \(\mathfrak{p}_3 = \mathfrak{p}_i'\) . (Argue by reductio ad absurdum: in the opposite case, there would exist \(y \notin \mathfrak{p}_3\) such that \(y \in \mathfrak{q}_i\) for all \(i \geqslant s + 1\) ; then \(x_2y \in \mathfrak{a}\) and prove that this relation contradicts \(y \notin \mathfrak{p}_3\) .)
\end{enumerate}

Conclude that \(\mathfrak{p}_3\) is a non-minimal element of \(\operatorname{Ass}_f(A / a)\) ; we denote by \(\mathfrak{q}_3'\) the primary component of \(a\) in \(A\) which corresponds to it in the above decomposition.

\begin{enumerate}
\def\labelenumi{(\alph{enumi})}
\setcounter{enumi}{2}
\tightlist
\item
  Let \(\mathfrak{b} = \mathfrak{p}_2 \cap \mathfrak{q}_3'\) ; on the other hand let \(m > 0\) be such that \(x_3^m \in \mathfrak{q}_3'\) and let \(\mathfrak{c} = \mathfrak{b} + Ax_3^{2m}\) . Show that the saturation of \(\mathfrak{c}\) with respect to \(\mathfrak{p}_3\) is a \(\mathfrak{p}_3\) -primary ideal \(\mathfrak{q}_3''\) (show that every element of \(\mathfrak{p}_3\) is a power in \(\mathfrak{q}_3''\) by using the fact that the saturation of \(\mathfrak{p}_2 + Ax_3\) with respect to \(\mathfrak{p}_3\) is primary). Prove on the other hand that \(x_3^m \notin \mathfrak{q}_3''\) and that \(\mathfrak{b} = \mathfrak{p}_2 \cap \mathfrak{q}_3''\) ; then there are two distinct reduced primary decompositions of the ideal \(\mathfrak{b}\) .
\end{enumerate}

\begin{enumerate}
\def\labelenumi{\arabic{enumi}.}
\setcounter{enumi}{24}
\tightlist
\item
  Let A be a Laskerian ring. Let \(a\) be an ideal of A and \(a = \bigcap_{1 \leqslant i \leqslant n} q_i\) a reduced primary decomposition of \(a\) in A; let \(p_i\) denote the radical of \(q_i\) ; suppose that the minimal elements of \(\operatorname{Ass}_f(A/a)\) are the \(p_i\) such that \(1 \leqslant i \leqslant s\) and that \(s < n\) . We write \(b = \bigcap_{1 \leqslant i \leqslant s} q_i\) , \(c = \bigcap_{s+1 \leqslant i \leqslant n} q_i\) .
\end{enumerate}

\begin{enumerate}
\def\labelenumi{(\alph{enumi})}
\item
  Show that there exists \(x \in c\) such that \(x \notin p_i\) for \(1 \leqslant i \leqslant s\) and that \(a = b \cap (a + Ax)\) .
\item
  Suppose for example that \(p_{s+1}\) is minimal in \(\operatorname{Ass}_{f}(\mathbf{A}/c)\) . Show that \(xp_{s+1} + a \neq Ax + a\) and \(a = b \cap (xp_{s+1} + a)\) ; moreover, \(p_{s+1}\) is minimal in \(\operatorname{Ass}_{f}(\mathbf{A}/b)\) , where \(d = xp_{s+1} + a\) . Show that the saturation of d with respect to \(p_{s+1}\) is different from \(q_{s+1}\) . (Note that this saturation cannot contain x.)
\end{enumerate}

¶ 26. Let A be a Laskerian ring in which every ideal ≠A has a unique reduced primary decomposition in A.

\begin{enumerate}
\def\labelenumi{(\alph{enumi})}
\item
  Show that for every ideal \(\mathfrak{a} \neq \mathbf{A}\) all the elements of \(\operatorname{Ass}_f(\mathbf{A} / \mathfrak{a})\) are minimal in this set (use Exercise 25 (b)). Let \(\{0\} = \bigcap_{1 \leqslant i \leqslant n} q_i\) be a reduced primary decomposition in \(\mathbf{A}\) and let \(p_i\) be the radical of \(q_i\) ( \(1 \leqslant i \leqslant n\) ). Show that every prime ideal of \(\mathbf{A}\) distinct from the \(p_i\) is maximal (use Exercise 24 (c)). If \(p\) is a prime ideal of \(\mathbf{A}\) , \(q\) an ideal contained in \(p\) and such that \(p / q\) is a nil ideal in \(\mathbf{A} / q\) , show that every ideal \(\mathfrak{a}\) such that \(q \subset \mathfrak{a} \subset p\) is \(p\) -primary. Deduce that if one of the \(p_i\) is not maximal, \(p_i\) is the only \(p_i\) -primary ideal and \(p_i^2 = p_i\) (use Exercise 16).
\item
  Show that in A every increasing sequence of ideals equal to their radicals is stationary. Deduce that every strictly increasing sequence \((\mathfrak{a}_n)\) of ideals such that every quotient \(\mathfrak{a}_n / \mathfrak{a}_{n-1}\) contains non-nilpotent elements in \(\mathbf{A} / \mathfrak{a}_{n-1}\) is finite. Conclude from this that for every prime ideal \(\mathfrak{p}\) of A there exists a primary ideal with radical \(\mathfrak{p}\) which is finitely generated. In particular, each of the \(\mathfrak{p}_i\) which is not maximal is finitely generated and therefore contains an idempotent \(e_i\) which generates it (Chapter II, § 4, Exercise 15).
\item
  Show that, for \(i \neq j \mathfrak{p}_i + \mathfrak{p}_j = \mathbf{A}\) (consider \(e_i + e_j - e_i e_j\) if \(\mathfrak{p}_i\) and \(\mathfrak{p}_j\) are not maximal). Conclude that \(\mathbf{A}\) is isomorphic to the product of a finite number of rings \(\mathbf{A}_i = \mathbf{A} / \mathfrak{q}_i\) such that: either \(\mathbf{A}_i\) is a local ring whose maximal ideal is a nil ideal; or \(\mathbf{A}\) is an integral domain in which every prime ideal \(\neq \{0\}\) is maximal and in which every ideal \(\neq\{0\}\) is only contained in a finite number of maximal ideals * (cf.~§ 1, Exercise 2). * Obtain the converse.
\end{enumerate}

¶ 27. Let M be an A-module and Q a p-primary submodule of M. Q is said to be of exponent m if m is the smallest integer k such that \(p^{k}M \subset Q\) (cf.~Exercise 3); if there exists no integer k with this property, Q is said to be of infinite exponent. If Q is of finite exponent, Q is also said to be strongly primary.

\begin{enumerate}
\def\labelenumi{(\alph{enumi})}
\tightlist
\item
  Show that, if \(\mathbf{Q} \colon \mathfrak{p} = \mathbf{Q}\) , \(\mathbf{Q}\) is of infinite exponent and that there exists \(\alpha \in \mathfrak{p}\) such that \(\mathbf{Q} \colon (\alpha)\) is distinct from \(\mathbf{Q}\) and is a \(\mathfrak{p}\) -primary submodule of \(\mathbf{M}\) of infinite exponent.
\end{enumerate}

* (b) Give an example of a ring A and a p-primary ideal q of infinite exponent and such that q: p = q (cf.~§ 1, Exercise 2). *

\begin{enumerate}
\def\labelenumi{(\alph{enumi})}
\setcounter{enumi}{2}
\item
  Let K be a field, A the polynomial ring \(K[X_{n}]_{n \in N}\) , p the prime ideal of A generated by the \(X_{n}\) and q the p-primary ideal generated by the products \(X_{i}X_{j} (i \neq j)\) and the elements \(X_{n}^{n+1}\) . Show that for all \(\alpha \in p - q\) , \(q: (\alpha)\) is of finite exponent, that \(q: p \neq q\) and that q: p is a p-primary ideal of infinite exponent.
\item
  Let \(\mathbf{Q}\) be a \(p\) -primary submodule of \(\mathbf{M}\) of finite exponent \(e\) . Show that the \(p^k\mathbf{M}\) for \(k < e\) are all distinct and that the submodules \(\mathbf{Q}: p^k\) for \(k < e\) are all distinct. For every ideal \(a\) of \(A\) show that either \(Q: a = Q\) or \(Q: a\) is a \(p\) -primary submodule in \(M\) distinct from \(M\) and of exponent \(\leqslant e - 1\) (observe that if \(a \subset p\) then \(p^{e-1}\mathbf{M} \subset Q: a\) ).
\end{enumerate}

¶ 28. A finitely generated A-module E is called strongly Laskerian if every submodule of E admits a primary decomposition in E all of whose elements are strongly primary (Exercise 27).

\begin{enumerate}
\def\labelenumi{(\alph{enumi})}
\tightlist
\item
  Show that for a finitely generated A-module E to be strongly Laskerian, it is necessary and sufficient that it satisfy axiom (LA \(_{II}\) ) of Exercise 23 and the axiom:
\end{enumerate}

\((\mathbf{L}\mathbf{A}_{\mathrm{III}})\) For every sequence \((\mathfrak{b}_k)\) of ideals of A and every submodule F of E, the increasing sequence of submodules \(\mathbf{F}:(\mathfrak{b}_1\mathfrak{b}_2\dots \mathfrak{b}_k)\) is stationary.

(To show that the conditions are necessary, use Exercise 27 (d). To see that they are sufficient, prove first that \((\mathrm{LA}_{\mathrm{III}})\) implies \((\mathrm{LA}_{\mathrm{I}})\) and then that it implies that every primary ideal is of finite exponent, using Exercise 27 (a).)

\begin{enumerate}
\def\labelenumi{(\alph{enumi})}
\setcounter{enumi}{1}
\item
  Show that the ring A defined in Exercise 23 (c) satisfies \((\mathbf{LA}_{\mathrm{III}})\) .
\item
  Let \(\mathbf{K}\) be a field, \(\mathbf{V}\) an infinite dimensional vector \(\mathbf{K}\) -space and \(\mathbf{A} = \mathbf{K} \oplus \mathbf{V}\) the ring whose multiplication is defined by
\end{enumerate}

\[
(a, x) (a ^ {\prime}, x ^ {\prime}) = (a a ^ {\prime}, a x ^ {\prime} + a ^ {\prime} x).
\]

Show that in A every ideal \(\neq A\) is strongly primary, although A is not Noetherian.

\begin{enumerate}
\def\labelenumi{(\alph{enumi})}
\setcounter{enumi}{3}
\item
  Let E be a finitely generated A-module and \(E' \subset E\) a finitely generated submodule. For E to be strongly Laskerian, it is necessary and sufficient that \(E'\) and \(E/E'\) be so. In particular, if A is a strongly Laskerian ring, every finitely generated A-module is strongly Laskerian. If E is a strongly Laskerian A-module, \(S^{-1}E\) is a strongly Laskerian \(S^{-1}A\) -module for every multiplicative subset S of A.
\item
  Generalize Exercises 3 and 4 to strongly Laskerian modules.
\end{enumerate}

¶ 29. Let A be a strongly Laskerian ring (Exercise 28).

\begin{enumerate}
\def\labelenumi{(\alph{enumi})}
\item
  Let \(\mathfrak{a}, \mathfrak{b}\) be two ideals of \(\mathbf{A}\) , \(\mathfrak{ab} = \bigcap_{i=1}^{n} q_i\) a reduced primary decomposition of \(\mathfrak{ab}\) in \(\mathbf{A}\) and \(\mathfrak{c}\) the intersection of those \(q_i\) whose radicals do not contain \(\mathfrak{b}\) ; show that \(\mathfrak{a} \subset \mathfrak{c}\) (if \(q_i\) is such that its radical \(p_i\) does not contain \(\mathfrak{b}\) and \(x \in \mathfrak{b}\) and \(x \notin p_i\) , consider the product \(ax\) ).
\item
  Deduce from (a) that there exist three ideals \(\mathfrak{a}', \mathfrak{b}', \mathfrak{c}\) and an integer \(s > 0\) such that \(\mathfrak{a}^s \subset \mathfrak{a}', \mathfrak{b}^s \subset \mathfrak{b}', (\mathfrak{a} + \mathfrak{b})^s \subset \mathfrak{c}\) and
\end{enumerate}

\[
a b = a ^ {\prime} \cap b = a \cap b ^ {\prime} = a \cap b \cap c.
\]

\begin{enumerate}
\def\labelenumi{(\alph{enumi})}
\setcounter{enumi}{2}
\item
  Show that for every finite family \((\mathfrak{a}_k)\) of ideals of A, there exists an integer \(m > 0\) such that \(\bigcap_{k} \mathfrak{a}_k^m \subset \prod_{k} \mathfrak{a}_k\) (apply (b) arguing by induction on the number of ideals \(\mathfrak{a}_k\) ).
\item
  For every ideal \(\mathfrak{a}\) of \(\mathbf{A}\) show that the intersection \(\mathfrak{b} = \bigcap_{n=1} \mathfrak{a}^n\) is the ideal of the \(x \in \mathbf{A}\) such that \(x \in x\mathfrak{a}\) (apply (b) to the product \((\mathbf{A}x)\mathfrak{a}\) ). Deduce that \(\mathfrak{b}\) is the intersection of the primary components of \(\{0\}\) (for a reduced primary decomposition of \(\{0\}\) ) whose radicals meet \(1 + \mathfrak{a}\) .
\item
  Deduce from (d) that, for every prime ideal \(\mathfrak{p}\) of \(\mathbf{A}\) , the intersection of the symbolic powers \(\mathfrak{p}^{(n)}\) (Exercise 18 (b)) for \(n\in\mathbf{N}\) is the ideal of those \(x\in\mathbf{A}\) such that \(x\in x\mathfrak{p}\) (consider the ring \(\mathbf{A}_{\mathfrak{p}}\) ).
\end{enumerate}

\begin{enumerate}
\def\labelenumi{\arabic{enumi}.}
\setcounter{enumi}{29}
\tightlist
\item
  Let M be an A-module and let N be a submodule of M admitting a reduced primary decomposition \(N = \bigcap_{i=1}^{n} Q_{i}\) ; let \(p_{i}\) be the prime ideal associated with \(Q_{i}\) .
\end{enumerate}

\begin{enumerate}
\def\labelenumi{(\alph{enumi})}
\item
  Show that, if \(\mathfrak{b}\) is an ideal of \(\mathbf{A}\) which is not contained in any of the \(\mathfrak{p}_i\) , then \(\mathbf{N}:\mathfrak{b} = \mathbf{N}\) (cf.~Exercise 17).
\item
  If each of the \(Q_{i}\) is strongly primary in M, show conversely that, if N: b = N, b is not contained in any of the \(p_{i}\) . (Note that, if b ⊂ \(p_{i}\) and P = \(\bigcap_{j \neq i} Q_{j}\) , then \(b^rP \subset N\) for some suitable integer r and deduce that, if N: b = N, then P = N.) Deduce that there then exists \(\beta \in b\) such that N: ( \(\beta\) ) = N.
\end{enumerate}

¶ 31. (a) Let A be a Noetherian ring, m a maximal ideal of A and p a prime ideal contained in m; show that, if the Hausdorff completion \(\hat{A}\) of A with respect to the m-adic topology is an integral domain, the filter base of symbolic power \(p^{(n)}\) of p tends to 0 with respect to the m-adic topology on A. (Reduce it to the case where \(\mathbf{A}\) is local; let \(q\) be a prime ideal of \(\hat{\mathbf{A}}\) whose trace on \(\mathbf{A}\) is \(p\) and for all \(n > 0\) let \(c_{n} \notin p\) be such that \(c_{n} p^{(n)} \subset p^{n}\) ; show that \(c_{n} p^{(n)} \hat{\mathbf{A}} \subset q^{(n)}\) . Use Exercise 29 (e) and Chapter III, § 2, no. 7, Proposition 8.)

\begin{enumerate}
\def\labelenumi{(\alph{enumi})}
\setcounter{enumi}{1}
\item
  Let A be a Noetherian ring, n an ideal of A, p a minimal prime ideal in \(\operatorname{Ass}(A/n)\) and \(m_{0}\) , m two maximal ideals of A containing p; let \(A(n)\) , \(A(m_{0})\) , \(A(m)\) denote the Hausdorff completions of A with respect to the n-adic, \(m_{0}\) -adic, m-adic topologies respectively. Let z be an element of \(A(n)\) whose canonical image in \(A(m_{0})\) is zero. Show that if \(A(m)\) is an integral domain, the canonical image of z in \(A(m)\) is also zero. (Take z as the limit of a sequence \((x_{n})\) of elements of A such that \(x_{n}-x_{m}\in n^{m}\) for \(n\geqslant m\) and \(x_{n}\in m_{0}^{n}\) ; deduce that \(x_{m}\) belongs to the closure of \(n^{m}\) with respect to the \(m_{0}\) -adic topology; using Exercise 8, conclude that \(x_{m}\in p^{(m)}\) for all m and complete with the aid of (a).)
\item
  Let A be a Zariski ring and r a defining ideal of A; suppose that \(\operatorname{Spec}(A)\) is connected and that for every maximal ideal m of A the Hausdorff completion \(A(m)\) of A with respect to the m-adic topology is an integral domain. Show then that the completion \(\hat{A}\) of A with respect to the r-adic topology is an integral domain. (Show that for every maximal ideal m of A, the canonical homomorphism \(\hat{A} \to A(m)\) is injective. For this, we may assume that r is the intersection of a finite number of prime ideals which are minimal in \(\operatorname{Ass}(A/r)\) , let \(r = p_{1} \cap p_{2} \cap \cdots \cap p_{s}\) ; show that the hypothesis on \(\operatorname{Spec}(A)\) implies for each i the existence of a maximal ideal \(m_{i}\) containing \((p_{1} \cap \cdots \cap p_{i-1}) + p_{i}\) . Using (b), show that, if the canonical image of \(z \in \hat{A}\) in \(A(m)\) is zero and for example \(p_{1} \subset m\) , then the canonical image of z in each of the \(A(m_{i})\) is zero, then conclude that the canonical image of z in \(A(m')\) is zero for every maximal ideal \(m'\) of A.)
\end{enumerate}

\begin{enumerate}
\def\labelenumi{\arabic{enumi}.}
\setcounter{enumi}{31}
\tightlist
\item
  A ring \(\mathbf{A}\) is called primary if it is a local ring and its maximal ideal \(\mathfrak{m}\) is a nil ideal (*) . Every ideal of \(\mathbf{A}\) contained in \(\mathfrak{m}\) is then \(\mathfrak{m}\) -primary in \(\mathbf{A}\) .
\end{enumerate}

\begin{enumerate}
\def\labelenumi{(\alph{enumi})}
\item
  Let A be a strongly Laskerian primary ring, so that in particular the maximal ideal m is nilpotent. Show that, if a, b are two ideals contained in m, necessarily a: b ≠ a (use Exercise 30 (b)).
\item
  A primary Noetherian ring is Artinian. Deduce that, for every finitely generated module M over a Noetherian ring A and every p-primary submodule N of M, there exists an integer m such that every strictly decreasing sequence \((\mathrm{M}_{i}/\mathrm{N})_{0\leqslant i\leqslant k}\) of submodules of M/N for which the \(M_{i}\) are p-primary has length \(k\leqslant m\) (consider the \(A_{p}\) -module \(M_{p}\) and note that \(M_{p}/N_{p}\) is annihilated by a power of \(pA_{p}\) ).
\item
  Show that an Artinian ring A in which every prime ideal \(\neq\{0\}\) is maximal is the direct composition of a finite number of primary rings.
\end{enumerate}

\begin{enumerate}
\def\labelenumi{\arabic{enumi}.}
\setcounter{enumi}{32}
\tightlist
\item
  Let A be a commutative ring. An ideal a of A is called primal if, in A/a, the set of divisors of 0 is an ideal p/a; the ideal p is then prime.
\end{enumerate}

\begin{enumerate}
\def\labelenumi{(\alph{enumi})}
\item
  Show that every primary ideal and every irreducible ideal is primal.
\item
  Let \(\mathfrak{a}\) be an ideal of \(\mathbf{A}\) such that the set of ideals \(\mathfrak{b}\) for which \(\mathfrak{a}: \mathfrak{b} \neq \mathfrak{a}\) has a greatest element \(\mathfrak{p}\) ; show that \(\mathfrak{p}\) is prime and that \(\mathfrak{a}\) is primal. * Give an example of a primal ideal \(\mathfrak{a}\) for which the above condition is not satisfied (cf.~§ 1, Exercise 2). *
\item
  In a Noetherian ring A characterize the primal ideals and give an example of a non-primary primal ideal (cf.~§ 1, Exercise 11 (b)).
\end{enumerate}

\begin{enumerate}
\def\labelenumi{\arabic{enumi}.}
\setcounter{enumi}{33}
\tightlist
\item
  In a commutative ring A an ideal a is called quasi-prime if for every ordered pair of ideals b, c of A, the relation \(b \cap c \subset a\) implies \(b \subset a\) or \(c \subset a\) ; every prime ideal is quasi-prime and every quasi-prime ideal is irreducible. If the Chinese remainder theorem is valid in A, show that every irreducible ideal is quasi-prime (Algebra, Chapter VI, §1, Exercise 25).
\end{enumerate}

* 35. Let K be a commutative field, B the valuation ring, whose order group is R, consisting of the ``formal power series'' \(\sum_{x} c_{x} T^{x}\) , where \(c_{x} \in \mathbf{K}\) , \(x \in \mathbf{R}_{+}\) and the set of \(x\) such that \(c_{x} \neq 0\) is well ordered; the ring A defined in §1, Exercise 2 is a subring of B and B is a flat A-module. If C is the A-module defined in §1, Exercise 2, show that \(\operatorname{Ass}_{\mathbf{B}}(\mathbf{C} \otimes_{\mathbf{A}} \mathbf{B}) \neq \varnothing\) although \(\operatorname{Ass}_{\mathbf{A}}(\mathbf{C}) = \varnothing\) (consider an element of C \(\otimes_{\mathbf{A}} \mathbf{B} = \mathbf{B}/\mathfrak{a}\mathbf{B}\) the canonical image of an element of B of valuation \(\alpha\) ). *

